\documentclass[12pt]{article}
\usepackage[margin=1in]{geometry}
\usepackage{enumitem}
\usepackage{titlesec}
\usepackage{parskip}
\usepackage{graphicx}
\usepackage{xcolor}
\usepackage{hyperref}
\usepackage{fontenc}

\title{\textbf{Diermeier and Vlaicu 2011 RG Notes}\\
\large Diermeier, Daniel, and Razvan Vlaicu. ``Parties, Coalitions, and the Internal Organization of Legislatures'' [2011].\\
\textit{American Political Science Review} 105(2): 359-380.}
\author{}
\date{}

\begin{document}

\maketitle

PPM = pivotal politics model: majoritarian systems, including parties in those systems, will reflect the preferences of the majority and hence the median legislator

Provides a link between pivotal politics, procedural cartel theory, and conditional party government theories

Existing theory: parties viewed as either floor voting coalitions, passing policy by means of disciplining party members, or procedural coalitions, facilitating a targeted legislative agenda. (Again considering the collective action problem in legislatures: an individual legislator could go off on their own to extract rents, but in doing so might sacrifice the good of the party as a whole)

Essentially, parties are not necessary / their functions are ``de facto'' performed when a group of sufficiently ideologically aligned legislators exists such that the value of the collective action defect is significantly lesser. This requires two conditions:

(1) Legislative activities are time-consuming and therefore costly

(2) Legislators are polarized in preferences even organically (that is, without party influence)

More broadly an argument of agenda-setting authority: time, capacity, and procedural constraints provide outsized influence for agenda-setters. Also a first-mover incentive, as amendments can be costly to amenders

Two-stage modeling approach: (1) model of preferences, where institutions are not influential, and (2) model of policy-making, where institutions are so. (Think John Krasinski's -- ``What you want, and how to get it.'' (A positivist critique: does it matter what I want if I'm not able to get it?)

Krehbiel: we can treat institutional arrangements as policy preferences themselves, essentially, and so their equilibrium outcomes will also represent the median legislative preferences

\end{document}
